Un marchand de glaces donne le choix à ses clients d'acheter la glace dans un cornet ou un pot.

Trois types de glaces sont proposés : sorbet, glace au lait et crème glacée.

Le tableau suivant donne les effectifs des choix de 200 clients ayant acheté une seule boule de glace.

Les réponses aux questions suivantes seront données sous forme de fractions non nécessairement simplifiées.

On interroge un client au hasard.

On admet que tous les clients ont la même probabilité d'être interrogés.

\medskip

\begin{tblr}{colspec = {X[m,c]X[m,c]X[m,c]X[m,c]X[m,c]},hline{1}={2-Z}{solid},hline{2-Z}={solid},vline{1}={2-Z}{solid},vline{2-Z}={solid}}
	               & Sorbet & Glace au lait & Crème glacée & \textbf{Total} \\
	Cornet         & 16     & 64            & 80           & 160            \\
	Pot            & 12     & 4             & 24           & 40             \\
	\textbf{Total} & 28     & 68            & 104          & 200
\end{tblr}


\begin{enumerate}
\item Donner la probabilité que ce client ait choisi un sorbet.

\item Donner la probabilité que ce client ait choisi un sorbet dans un pot.

\item Sachant que le client interrogé a choisi un cornet, quelle est la probabilité qu'il ait choisi une glace au lait ?

\item Sachant que le client interrogé a choisi une crème glacée, quelle est la probabilité qu'il ait choisi un pot ?
\end{enumerate}
